\chapter{Complexity — P Subset P Poly Config Step}
%%% generated by the blueprint pipeline from TCSlib.Complexity.CircuitComplexity.PSubsetPPolyConfigStep
\section{Overview}

Declarations of the Lean module \texttt{TCSlib.Complexity.CircuitComplexity.PSubsetPPolyConfigStep}.

\section{Declarations}

\begin{definition}[Encoding an optional bit]
\lean{Complexity.CfgTableau.optBits}
\label{Complexity.CfgTableau.optBits}
\leanok
For an optional Boolean value $o$, define its two-bit encoding to be the list whose
first bit records whether $o$ is present and whose second bit is its value when present,
taking the second bit to be false when $o$ is absent.
\end{definition}

\begin{definition}[Constructing an optional bit]
\lean{Complexity.CfgTableau.optOf}
\label{Complexity.CfgTableau.optOf}
\leanok
Given Boolean values $s$ and $v$, define an optional Boolean value to be present with
value $v$ when $s$ is true, and absent when $s$ is false.
\end{definition}

\begin{theorem}[Reconstruction of an optional Boolean]
\lean{Complexity.CfgTableau.optOf_isSome_getD}
\label{Complexity.CfgTableau.optOf_isSome_getD}
\leanok
\uses{Complexity.CfgTableau.optOf}
\difficulty{4}
For every optional Boolean $o$, the optional Boolean formed by using whether $o$ is
present as the presence flag and its value (or false when absent) as the Boolean value
equals $o$.
\end{theorem}

\begin{theorem}[Optional Boolean from false flag]
\lean{Complexity.CfgTableau.optOf_false}
\label{Complexity.CfgTableau.optOf_false}
\leanok
\uses{Complexity.CfgTableau.optOf}
\difficulty{1}
For every Boolean $v$, choosing the false flag in the construction that returns `some v`
when the flag is true and `none` otherwise gives `none`.
\end{theorem}

\begin{definition}[Configuration kind]
\lean{Complexity.CfgKind}
\label{Complexity.CfgKind}
\leanok
For a natural number $k$, a configuration kind is either absent, present with no tape
index, or present with a tape index from a set of $k$ indices.
\end{definition}

\begin{definition}[Work-tape head update]
\lean{Complexity.cellHead}
\label{Complexity.cellHead}
\leanok
\uses{Complexity.Snapshot}
Let $M$ be a finite Turing machine over the Boolean alphabet, let $s$ be a snapshot of
$M$, let $\tau$ be one of its work tapes, and let $h_\ell,h_m,h_r$ be Boolean indicators
for the neighboring cells’ head positions. If $s$ is halted, the updated head indicator
is $h_m$; otherwise it is the indicator among $h_\ell,h_m,h_r$ selected by the
transition’s movement on tape $\tau$ (left selects $h_\ell$, stationary selects $h_m$,
and right selects $h_r$).
\end{definition}

\begin{definition}[Work-tape cell update]
\lean{Complexity.cellContent}
\label{Complexity.cellContent}
\leanok
\uses{Complexity.Snapshot}
Let $M$ be a finite Turing machine over the Boolean alphabet, let $s$ be a snapshot of
$M$, let $\tau$ be one of its work tapes, let $h_m$ indicate whether its head is at the
cell, and let $c$ be the cell’s optional Boolean content. If $s$ is halted or $h_m$ is
false, the updated content is $c$. Otherwise it is the value written by the transition
on tape $\tau$, if a value is specified, and is $c$ if no value is specified.
\end{definition}

\begin{definition}[Input-head update]
\lean{Complexity.inHead}
\label{Complexity.inHead}
\leanok
\uses{Complexity.Snapshot}
Let $M$ be a finite Turing machine over the Boolean alphabet, let $s$ be a snapshot of
$M$, and let $i_L,i_R,h_\ell,h_m,h_r$ be Boolean indicators, with $i_L$ and $i_R$
recording whether the cell is at the left and right input boundaries. If $s$ is halted,
the updated head indicator is $h_m$. Otherwise, a rightward transition gives $h_\ell$
or, at the right boundary, $h_m$; a stationary transition gives $h_m$; and a leftward
transition gives $h_r$ or, at the left boundary, $h_m$.
\end{definition}

\begin{definition}[Configuration tableau update rule]
\lean{Complexity.cfgF}
\label{Complexity.cfgF}
\leanok
\uses{Complexity.CfgKind, Complexity.CfgTableau.optBits, Complexity.CfgTableau.optOf, Complexity.Snapshot, Complexity.cellContent, Complexity.cellHead, Complexity.emitted, Complexity.inHead, Complexity.snapDecode, Complexity.snapEncode, Complexity.snapWidth, Complexity.stepState}
Let $M$ be a finite Turing machine over the Boolean alphabet, with $k$ work tapes, and
let $l$ be a finite Boolean list. The function indexed by configuration kinds returns a
Boolean list as follows. For a work-tape kind specifying tape $	au$, read the preceding
snapshot from the block of $l$ beginning at position $13$, of length equal to the number
of possible snapshots; determine the new head bit from positions $0$–$1$ when the marker
at position $0$ is set, and otherwise from the preceding snapshot and positions $2,5,8$.
Determine the cell content from the preceding snapshot and positions $5$–$7$, unless the
marker is set, in which case it is blank; the chain value is this content when the new
head bit is true, and is encoded by positions $11$–$12$ otherwise. The output is the
head bit followed by the two-bit encodings of the content and chain value. For the
input-tape kind, read the preceding snapshot from the block beginning at position $11$;
determine the new head bit from positions $0$–$1$ when the marker is set, and otherwise
from the preceding snapshot and positions $2,3,6,7,8$. The chain value is encoded by
positions $4$–$5$ when the new head bit is true and by positions $9$–$10$ otherwise, and
the output is the head bit followed by its two-bit encoding. For the absent kind, read
the preceding snapshot from the block beginning at position $4+2k$ and form a new
snapshot whose state is the initial state when position $0$ is true and otherwise the
successor state of the preceding snapshot, whose input symbol is encoded by positions
$2$–$3$, and whose work-tape symbols are encoded by positions $4+2	au$–$5+2	au$. The
output is the encoding of that snapshot followed by a bit that is false when position
$0$ is true and otherwise is the disjunction of position $1$ and whether the preceding
snapshot emits true. At any position beyond the end of $l$, the bit read is false.
\end{definition}

\begin{definition}[Machine configuration at time t]
\lean{Complexity.CfgTableau.cfgAt}
\label{Complexity.CfgTableau.cfgAt}
\leanok
\uses{Turing.Cfg, Turing.MultiTapeTM.initCfg, Turing.MultiTapeTM.runFrom}
Let $M$ be a deterministic Turing machine with Boolean input alphabet, let $w$ be a
finite Boolean input string, and let $t$ be a nonnegative integer. The configuration of
$M$ on input $w$ after $t$ steps is the machine's run from its initial configuration for
exactly $t$ steps.
\end{definition}

\begin{definition}[Work tape chain value]
\lean{Complexity.workChain}
\label{Complexity.workChain}
\leanok
\uses{Complexity.CfgTableau.cfgAt, Turing.Cfg.workTapeSymbols}
Let $M$ be a deterministic Turing machine with Boolean input alphabet, let $w$ be a
finite Boolean input string, and let $t$ be a nonnegative integer. For a work tape $	au$
and cell position $z\in\mathbb Z$, the work chain value is the symbol at cell $z$ in the
configuration after $t$ steps if that tape's head position is at most $z$, and is absent
otherwise.
\end{definition}

\begin{definition}[Work tape cell encoding]
\lean{Complexity.cellExp}
\label{Complexity.cellExp}
\leanok
\uses{Complexity.CfgTableau.cfgAt, Complexity.CfgTableau.optBits, Complexity.workChain}
Let $M$ be a deterministic Turing machine with Boolean input alphabet, let $w$ be a
finite Boolean input string, and let $T,t,r$ be nonnegative integers. For a work tape
$\tau$, the cell encoding at position $r-T$ consists of the Boolean indicator that the
tape head is at that position, followed by the Boolean encodings of the tape symbol
there and the work chain value there.
\end{definition}

\begin{definition}[Input tape chain value]
\lean{Complexity.inpChain}
\label{Complexity.inpChain}
\leanok
\uses{Complexity.CfgTableau.cfgAt, Turing.Cfg.inputSymbol}
Let $M$ be a deterministic Turing machine with Boolean input alphabet, let $w$ be a
finite Boolean input string, and let $t,p$ be nonnegative integers. The input chain
value at position $p$ is the input symbol in the configuration after $t$ steps if the
input head position is at most $p$, and is absent otherwise.
\end{definition}

\begin{definition}[Input tape position encoding]
\lean{Complexity.inpExp}
\label{Complexity.inpExp}
\leanok
\uses{Complexity.CfgTableau.cfgAt, Complexity.CfgTableau.optBits, Complexity.inpChain}
Let $M$ be a deterministic Turing machine with Boolean input alphabet, let $w$ be a
finite Boolean input string, and let $t,p$ be nonnegative integers. The input position
encoding at $p$ consists of the Boolean indicator that the input head is at position
$p$, followed by the Boolean encoding of the input chain value at $p$.
\end{definition}

\begin{definition}[Machine snapshot encoding]
\lean{Complexity.snapExp}
\label{Complexity.snapExp}
\leanok
\uses{Complexity.snapEncode, Complexity.snapshotAt, Complexity.tableauAccepted}
Let $M$ be a deterministic Turing machine with Boolean input alphabet, let $w$ be a
finite Boolean input string, and let $t$ be a nonnegative integer. The snapshot encoding
at time $t$ is the encoding of the machine snapshot at that time, with the Boolean
indicator of whether the tableau is accepted appended.
\end{definition}

\begin{theorem}[Bound on work tape head displacement]
\lean{Complexity.abs_workTapePos_le}
\label{Complexity.abs_workTapePos_le}
\leanok
\uses{Complexity.CfgTableau.cfgAt, Turing.Cfg.init, Turing.MultiTapeTM.runFrom_succ_eq_step', Turing.MultiTapeTM.runFrom_zero, Turing.MultiTapeTM.workTapePos_step_le}
\difficulty{4}
Let $M$ be a finite Turing machine over the binary alphabet, let $w$ be a finite binary
string, let $t\in\mathbb{N}$, and let $\tau$ be one of its work tapes. After running $M$
on $w$ for $t$ steps, the absolute value of the head position on tape $\tau$ is at most
$t$.
\end{theorem}

\begin{theorem}[Configuration Input Read as a Boundary-Aware Bit]
\lean{Complexity.inputSymbol_eq_inputBitAt}
\label{Complexity.inputSymbol_eq_inputBitAt}
\leanok
\uses{Complexity.inputBitAt, Turing.Cfg, Turing.Cfg.inputSymbol, Turing.FinTM.inputSymbol_at}
\difficulty{4}
Let $M$ be a finite-state-set Turing machine over the binary alphabet, let $w$ be a
binary input, and let $c$ be a configuration of $M$ on $w$. The symbol read by the input
head equals the boundary-aware input symbol of $w$ at the input-head position: it is
absent at position $0$ or beyond the final input position, and otherwise is the bit at
index one less than that position.
\end{theorem}

\begin{theorem}[One-step work-tape cell update]
\lean{Complexity.cell_step}
\label{Complexity.cell_step}
\leanok
\uses{Complexity.cellContent, Complexity.cellHead, Turing.Action.apply, Turing.Action.apply_workTapes, Turing.Cfg, Turing.Cfg.inputSymbol, Turing.Cfg.workTapeSymbols, Turing.MultiTapeTM.step}
\difficulty{5}
Let $M$ be a finite Turing machine over the binary alphabet, let $w$ be a finite bit
string, and let $c$ be a configuration of $M$ on input $w$. For each work tape $\tau$
and integer cell position $z$, the indicator that the head on tape $\tau$ is at $z$
after one step equals the cell-head rule applied to the current snapshot and the
indicators that the head was at $z-1$, $z$, and $z+1$. Also, the content of cell $z$ on
tape $\tau$ after one step equals the cell-content rule applied to the current snapshot,
the indicator that the head was at $z$, and the cell's current content.
\end{theorem}

\begin{theorem}[Head position after a positive move]
\lean{Complexity.moveInputPos_pos_val}
\label{Complexity.moveInputPos_pos_val}
\leanok
\uses{Turing.moveInputPos}
\difficulty{4}
Let $n\in\mathbb{N}$ and let $q$ be one of the $n+2$ valid head positions. After moving
in the positive direction, the resulting position, viewed as a natural number, is
$\min(q+1,n+1)$.
\end{theorem}

\begin{theorem}[Head position after a negative move]
\lean{Complexity.moveInputPos_neg_val}
\label{Complexity.moveInputPos_neg_val}
\leanok
\uses{Turing.moveInputPos}
\difficulty{4}
Let $n\in\mathbb{N}$ and let $q$ be one of the $n+2$ valid head positions. After moving
in the negative direction, the resulting position, viewed as a natural number, is $q-1$,
with natural-number subtraction.
\end{theorem}

\begin{theorem}[Input Head Position After One Step]
\lean{Complexity.inp_step}
\label{Complexity.inp_step}
\leanok
\uses{Complexity.inHead, Complexity.moveInputPos_neg_val, Complexity.moveInputPos_pos_val, Turing.Action.apply, Turing.Cfg, Turing.Cfg.inputSymbol, Turing.Cfg.workTapeSymbols, Turing.MultiTapeTM.step, Turing.moveInputPos_zero}
\difficulty{5}
Let $M$ be a binary-input machine, let $w$ be an input word, and let $c$ be a
configuration on $w$. For $p\le |w|+1$, the indicator that the input-head position after
one step is $p$ equals the following Boolean condition: if $c$ is halted, the current
position is $p$; otherwise, inspect the transition's input-head move. A right move
satisfies the condition when either $p=0$ or the current position is $p-1$ with $p\ge1$;
a stationary move satisfies it when the current position is $p$; and a left move
satisfies it when either $p=|w|+1$ or the current position is $p+1$.
\end{theorem}

\begin{theorem}[Acceptance recurrence]
\lean{Complexity.tableauAccepted_succ}
\label{Complexity.tableauAccepted_succ}
\leanok
\uses{Complexity.emitted, Complexity.snapshotAt, Complexity.tableauAccepted}
\difficulty{4}
Let $M$ be a finite-state Turing machine over the binary alphabet, let $w$ be a binary
input word, and let $t\in\mathbb{N}$. The machine has emitted the bit $1$ at some time
strictly before $t+1$ if and only if it has emitted $1$ at some time strictly before
$t$, or its transition at time $t$ emits $1$.
\end{theorem}

\begin{theorem}[Work-chain recurrence]
\lean{Complexity.workChain_eq}
\label{Complexity.workChain_eq}
\leanok
\uses{Complexity.CfgTableau.cfgAt, Complexity.workChain, Turing.Cfg.workTapeSymbols}
\difficulty{4}
Let $M$ be a finite-state Turing machine with Boolean work-tape symbols, let $w$ be a
finite Boolean input, let $t\in\mathbb{N}$, let $\tau$ be a work-tape index, and let
$z\in\mathbb{Z}$. The work-chain value at $z$ after $t$ steps equals the contents of
cell $z$ on tape $\tau$ if that tape's head is at $z$; otherwise it equals the
work-chain value at $z-1$.
\end{theorem}

\begin{theorem}[Recurrence for the input chain]
\lean{Complexity.inpChain_eq}
\label{Complexity.inpChain_eq}
\leanok
\uses{Complexity.CfgTableau.cfgAt, Complexity.inpChain, Complexity.inputBitAt, Complexity.inputSymbol_eq_inputBitAt}
\difficulty{5}
Let $M$ be a finite-state multi-tape Turing machine over the binary alphabet, and let
$w$ be a binary input, with $t,p\in\mathbb N$. The input-chain value at position $p$
after $t$ steps equals the input symbol at position $p$ if the input head is at position
$p$; otherwise it is empty when $p=0$, and for $p>0$ it equals the input-chain value at
position $p-1$.
\end{theorem}

\begin{theorem}[Initial work-cell update rule]
\lean{Complexity.cfgF_cell_zero}
\label{Complexity.cfgF_cell_zero}
\leanok
\uses{Complexity.CfgTableau.cfgAt, Complexity.CfgTableau.optBits, Complexity.CfgTableau.optOf_false, Complexity.cfgF, Complexity.workChain, Turing.Cfg.init, Turing.Cfg.workTapeSymbols, Turing.MultiTapeTM.runFrom_zero}
\difficulty{3}
Let $M$ be a finite-state binary Turing machine with $k$ work tapes, a binary input $w$,
a work tape $\tau$, a cell position $z\in\mathbb Z$, and a finite list of bits $l$. If
the first bit of $l$ is true, its second bit records whether $z=0$, and its designated
chain bit is false, then the local work-cell update rule applied to $l$ yields the
time-zero encoding of whether the work-tape head is at $z$, the content of cell $z$, and
the work-chain value at $z$.
\end{theorem}

\begin{theorem}[Work-cell update rule]
\lean{Complexity.cfgF_cell_succ}
\label{Complexity.cfgF_cell_succ}
\leanok
\uses{Complexity.CfgTableau.cfgAt, Complexity.CfgTableau.optBits, Complexity.CfgTableau.optOf_isSome_getD, Complexity.cell_step, Complexity.cfgF, Complexity.snapDecode, Complexity.snapWidth, Complexity.snapshotAt, Complexity.workChain, Complexity.workChain_eq, Turing.Cfg.inputSymbol, Turing.Cfg.workTapeSymbols, Turing.MultiTapeTM.runFrom_succ_eq_step', Turing.MultiTapeTM.step}
\difficulty{5}
Let $M$ be a finite-state binary Turing machine with $k$ work tapes, and let $w$ be a
binary input. Fix a time $t\in\mathbb N$, a work tape $\tau$, a cell position
$z\in\mathbb Z$, and a finite list of bits $l$. Suppose the first bit of $l$ is false,
its snapshot field decodes to the machine's snapshot on $w$ at time $t$, and its
designated neighboring-head bits record whether the tape head at time $t$ is at $z-1$,
$z$, and $z+1$. Suppose also that the designated content bits record whether cell $z$
contains a bit and, if so, its value; and that the designated chain bits record whether
the time-$t+1$ work-chain value at $z-1$ exists and, if so, its value. Then the local
work-cell update rule applied to $l$ yields the encoding of the time-$t+1$ head
indicator, cell content, and work-chain value at $z$.
\end{theorem}

\begin{theorem}[Initial input-position update rule]
\lean{Complexity.cfgF_inp_zero}
\label{Complexity.cfgF_inp_zero}
\leanok
\uses{Complexity.CfgTableau.cfgAt, Complexity.CfgTableau.optBits, Complexity.CfgTableau.optOf_isSome_getD, Complexity.cfgF, Complexity.inpChain, Complexity.inpChain_eq, Complexity.inputBitAt, Turing.Cfg.init, Turing.MultiTapeTM.runFrom_zero}
\difficulty{4}
Let $M$ be a finite-state binary Turing machine, let $w$ be a binary input, and let
$p\in\mathbb N$ and a finite list of bits $l$ be given. If the first bit of $l$ is true,
its second bit records whether $p=1$, its input-symbol bits encode the symbol at
position $p$ in $w$, and its designated chain bits encode the presence and value of the
time-zero input-chain value at $p-1$ when $p>0$ (and are absent when $p=0$), then the
local input-position update rule applied to $l$ yields the encoding of whether the
time-zero input head is at $p$ and of the time-zero input-chain value at $p$.
\end{theorem}

\begin{theorem}[Input-position update rule]
\lean{Complexity.cfgF_inp_succ}
\label{Complexity.cfgF_inp_succ}
\leanok
\uses{Complexity.CfgTableau.cfgAt, Complexity.CfgTableau.optBits, Complexity.CfgTableau.optOf_isSome_getD, Complexity.cfgF, Complexity.inpChain, Complexity.inpChain_eq, Complexity.inp_step, Complexity.inputBitAt, Complexity.snapDecode, Complexity.snapWidth, Complexity.snapshotAt, Turing.Cfg.inputSymbol, Turing.Cfg.workTapeSymbols, Turing.MultiTapeTM.runFrom_succ_eq_step', Turing.MultiTapeTM.step}
\difficulty{4}
Let $M$ be a finite-state binary Turing machine with $k$ work tapes, let $w$ be a binary
input, and fix $t,p\in\mathbb N$ with $p\le |w|+1$ and a finite list of bits $l$.
Suppose the first bit of $l$ is false, its snapshot field decodes to the machine's
snapshot on $w$ at time $t$, its boundary bits indicate whether $p=0$ and $p=|w|+1$, and
its input-symbol bits encode the symbol read from $w$ at position $p$. Suppose its
neighboring-head bits encode whether the input head at time $t$ is at $p-1$, at $p$, or
at $p+1$ (with the left-neighbor condition also requiring $p\ge1$). Finally, suppose its
designated chain bits encode the presence and value of the time-$(t+1)$ input-chain
value at $p-1$, using no value when $p=0$. Then the local input-position update rule
applied to $l$ yields the encoding of whether the input head at time $t+1$ is at $p$ and
of the time-$(t+1)$ input-chain value at $p$.
\end{theorem}

\begin{theorem}[Snapshot update rule]
\lean{Complexity.cfgF_snap}
\label{Complexity.cfgF_snap}
\leanok
\uses{Complexity.CfgTableau.cfgAt, Complexity.CfgTableau.optOf_isSome_getD, Complexity.cfgF, Complexity.snapDecode, Complexity.snapExp, Complexity.snapWidth, Complexity.snapshotAt, Complexity.snapshotAt_state_succ, Complexity.tableauAccepted, Complexity.tableauAccepted_succ, Turing.Cfg.inputSymbol, Turing.Cfg.workTapeSymbols}
\difficulty{4}
Let $M$ be a finite-state binary Turing machine with $k$ work tapes, let $w$ be a binary
input, and fix $t\in\mathbb N$ and a finite list of bits $l$. Suppose the first bit
records whether $t=0$; when $t\ne0$, the second bit records whether the tableau had
already accepted by time $t-1$. Suppose the remaining designated bits encode the input
symbol and each work-tape symbol in the configuration at time $t$, by presence and
value, and, when $t\ne0$, the snapshot field decodes to the machine's snapshot at time
$t-1$. Then the snapshot update rule applied to $l$ yields the snapshot encoding at time
$t$ followed by the acceptance indicator at time $t$.
\end{theorem}
